\documentclass{article}

% if you need to pass options to natbib, use, e.g.:
%     \PassOptionsToPackage{numbers, compress}{natbib}
% before loading neurips_2025

% The authors should use one of these tracks.
% Before accepting by the NeurIPS conference, select one of the options below.
% 0. "default" for submission
 \usepackage{neurips_2025}
% the "default" option is equal to the "main" option, which is used for the Main Track with double-blind reviewing.
% 1. "main" option is used for the Main Track
%  \usepackage[main]{neurips_2025}
% 2. "position" option is used for the Position Paper Track
%  \usepackage[position]{neurips_2025}
% 3. "dandb" option is used for the Datasets & Benchmarks Track
 % \usepackage[dandb]{neurips_2025}
% 4. "creativeai" option is used for the Creative AI Track
%  \usepackage[creativeai]{neurips_2025}
% 5. "sglblindworkshop" option is used for the Workshop with single-blind reviewing
 % \usepackage[sglblindworkshop]{neurips_2025}
% 6. "dblblindworkshop" option is used for the Workshop with double-blind reviewing
%  \usepackage[dblblindworkshop]{neurips_2025}

% After being accepted, the authors should add "final" behind the track to compile a camera-ready version.
% 1. Main Track
 % \usepackage[main, final]{neurips_2025}
% 2. Position Paper Track
%  \usepackage[position, final]{neurips_2025}
% 3. Datasets & Benchmarks Track
 % \usepackage[dandb, final]{neurips_2025}
% 4. Creative AI Track
%  \usepackage[creativeai, final]{neurips_2025}
% 5. Workshop with single-blind reviewing
%  \usepackage[sglblindworkshop, final]{neurips_2025}
% 6. Workshop with double-blind reviewing
%  \usepackage[dblblindworkshop, final]{neurips_2025}
% Note. For the workshop paper template, both \title{} and \workshoptitle{} are required, with the former indicating the paper title shown in the title and the latter indicating the workshop title displayed in the footnote.
% For workshops (5., 6.), the authors should add the name of the workshop, "\workshoptitle" command is used to set the workshop title.
% \workshoptitle{WORKSHOP TITLE}

% "preprint" option is used for arXiv or other preprint submissions
 % \usepackage[preprint]{neurips_2025}

% to avoid loading the natbib package, add option nonatbib:
%    \usepackage[nonatbib]{neurips_2025}

\usepackage[utf8]{inputenc} % allow utf-8 input
\usepackage[T1]{fontenc}    % use 8-bit T1 fonts
\usepackage{hyperref}       % hyperlinks
\usepackage{url}            % simple URL typesetting
\usepackage{booktabs}       % professional-quality tables
\usepackage{amsfonts}       % blackboard math symbols
\usepackage{amssymb}
\usepackage{nicefrac}       % compact symbols for 1/2, etc.
\usepackage{microtype}      % microtypography
\usepackage{xcolor}         % colors
\usepackage{amsmath}
\usepackage{amsthm}
\usepackage{graphicx}
\usepackage{array}
\usepackage{tabularx}
\makeatletter
\setlength{\@neuripsbelowcaptionskip}{10\p@}
\makeatother

% Note. For the workshop paper template, both \title{} and \workshoptitle{} are required, with the former indicating the paper title shown in the title and the latter indicating the workshop title displayed in the footnote.
\newtheorem{lemma}{Lemma}
\newtheorem{proposition}{Proposition}
\newtheorem{corollary}{Corollary}
\newcommand{\tv}{\mathrm{TV}}

\title{ADSD: Adversarial Deceleration of Speculative Decoding via Early-Rejection Prompt Optimization}


% The \author macro works with any number of authors. There are two commands
% used to separate the names and addresses of multiple authors: \And and \AND.
%
% Using \And between authors leaves it to LaTeX to determine where to break the
% lines. Using \AND forces a line break at that point. So, if LaTeX puts 3 of 4
% authors names on the first line, and the last on the second line, try using
% \AND instead of \And before the third author name.


\author{%
  Anonymous Authors \\
  Anonymous Institution \\
  Anonymous Address \\
  \texttt{anonymous@example.com} \\
}


\begin{document}


\maketitle


\begin{abstract}
Speculative decoding speeds up large language model inference only when a draft
model stays aligned with the target verifier. We study the opposite regime:
prompt-side attacks that deliberately collapse speculative acceptance and
thereby erase the latency benefit of lossless decoding. Our starting point is
the token-wise acceptance rule
$\alpha_i=\min(1,p(y_i\mid h_i)/q(y_i\mid h_i))$. We show that for a fixed
realized draft path, the expected accepted prefix length equals
$\sum_{k=1}^{\gamma}\prod_{i=1}^{k}\alpha_i$, and that under a fixed attack
budget it is always optimal to move rejection pressure to earlier draft
positions. This yields a simple mathematical principle: the most damaging bad
token is the earliest possible one. We further derive a smooth one-step
surrogate in which expected acceptance equals the overlap
$1-\tv(p,q)$ between draft and target distributions, leading to a
gradient-based relaxed prompt optimization method. In a local white-box setup
with \texttt{Qwen2.5-0.5B-Instruct-GPTQ-Int8} drafting for
\texttt{Qwen2.5-14B-Instruct-GPTQ-Int8}, a benign GSM8K-style prompt has
first-token acceptance $\alpha_1=1.0$, while the best adversarial seed we found
so far drops to $\alpha_1=0.160$, causing an $84\%$ first-token rejection
probability. On the saved 263-sample speculative GSM8K run, the ordinary
aggregate drafted-token acceptance rate is $49.66\%$. These results identify a
new prompt-side systems vulnerability for speculative decoding and provide the
first mathematical foundation for gradient-based early-rejection attacks.
\end{abstract}

\section{Introduction}

Speculative decoding accelerates large language model (LLM) inference by
letting a small draft model propose multiple tokens that are then verified by a
larger target model in parallel \citep{leviathan2023speculative}. The resulting
speedup is lossless at the output layer, but it is not free: it depends
critically on the degree of alignment between the draft distribution and the
target verifier. When draft-target agreement is high, long blocks of drafted
tokens are accepted and throughput improves substantially. When agreement
collapses, the verifier commits fewer drafted tokens, target-side correction
work increases, and the latency benefits of speculative decoding shrink.

Most prior work treats this acceptance bottleneck as an optimization problem to
be fixed. Recent systems improve the draft, adapt draft length, redesign
verification, or otherwise try to raise acceptance rates
\citep{liu2023online,goel2024direct,agrawal2024adaedl,bachmann2025judge,wang2025alignment,zhou2026hsd}.
This paper studies the complementary adversarial question: can an attacker
exploit the same acceptance bottleneck from the input side? In particular, can
a short prompt prefix preserve user-visible task behavior while forcing the
verifier into a substantially less efficient regime?

This attack surface is natural in deployment. The attacker does not need model
weights, scheduler privileges, or a custom decoding stack at inference time.
Instead, the attacker controls only the prompt text. Because speculative
decoding ties latency directly to the relation between the draft distribution
$q$ and the target distribution $p$, prompt-side inputs can act as
systems-level levers: they can induce overconfident draft proposals that the
target model repeatedly rejects, thereby converting speculative acceleration
into additional verifier work.

The technical obstacle is discrete. Gradient-based prompt optimization operates
most naturally over continuous token relaxations, but deployment requires a
discrete projected prefix. The main question is therefore not simply whether a
smooth surrogate can be written down, but when optimizing that surrogate yields
a strong \emph{discrete} attack after projection. This continuous-to-discrete
gap is the central methodological challenge for prompt-side attacks on
speculative decoding.

Our starting point is the token-wise verifier acceptance rule
\[
\alpha_i = \min\!\left(1, \frac{p(y_i \mid h_i)}{q(y_i \mid h_i)}\right),
\]
where $h_i$ is the speculative history and $y_i$ is the drafted token at step
$i$. Under draft sampling, the one-step expected acceptance satisfies
\[
\mathbb{E}[\alpha_i] = 1 - \tv(p_i, q_i),
\]
which immediately motivates an adversarial objective: maximize early
draft-target disagreement, since early verifier rejection is the most damaging
to end-to-end speculative efficiency. We formalize this intuition with a
continuous objective over relaxed prefixes, show that unconditional projection
guarantees are impossible via a prompt-waywardness proposition, and prove a
relaxation-gap bound under a TV-Lipschitz assumption that directly motivates a
projection-distortion penalty in the implemented optimizer.

On the algorithmic side, we instantiate this theory in two stages. The first
stage is a verifier-aligned soft-prefix optimizer that targets first-token
collapse at the assistant boundary using a straight-through relaxed prefix, a
soft-collapse term, TV, reverse KL, entropy sharpening, and an explicit
projection-distortion penalty. The second stage strengthens the attack into a
calibration-batch universal adversarial prefix optimized over target-generated
multi-timestamp rollouts. This universal-prefix pipeline combines a weighted
trajectory-level objective with a gradient-guided discrete seed-search stage,
allowing the optimizer to discover stronger discrete initializations before
continuous refinement.

We evaluate the attack in the reproduced HSD workspace using
\texttt{Qwen2.5-0.5B-Instruct-GPTQ-Int8} as the draft model and
\texttt{Qwen2.5-14B-Instruct-GPTQ-Int8} as the target model. The initial
Palmetto soft-prefix search consistently recovers the same five-token discrete
boundary attack across a sweep of projection-penalty strengths, reducing exact
first-token acceptance from $\alpha_1 = 0.1597$ for the manual comparison seed
to $\alpha_1 = 0.0288$ on the optimized search context. On a matched
263-sample token-wise GSM8K benchmark, that boundary attack raises mean sample
time from $23.10$s to $29.81$s while reducing aggregate acceptance from
$0.597$ to $0.509$ and block efficiency from $5.832$ to $4.961$.

The stronger result comes from the calibration-batch universal-prefix attack.
The best accuracy-preserving universal prefix raises mean sample time to
$38.61$s, lowers aggregate acceptance to $0.390$, lowers block efficiency to
$3.731$, and reduces decoding speed to $51.68$, while matching the benign
token-wise GSM8K accuracy exactly at $222/263$. In other words, a short
prompt-side prefix can preserve task accuracy while extracting a large
inference-cost penalty from a lossless speculative-decoding stack.

Our contributions are:
\begin{itemize}
\item We formulate prompt-side speculative-decoding slowdown as a
continuous-to-discrete adversarial optimization problem over verifier
acceptance rather than output semantics.
\item We prove a prompt-waywardness impossibility result and a relaxation-gap
bound that makes the projection distortion term in the optimizer theoretically
explicit.
\item We implement a verifier-aligned soft-prefix attack and a
calibration-batch universal-prefix extension with gradient-guided discrete
seed search, both of which are faithful to the released code path.
\item We show on a reproduced HSD/Qwen stack that prompt-only attacks can
materially degrade speculative efficiency, and that the strongest
accuracy-preserving universal prefix raises end-to-end runtime by $67.1\%$ on a
matched token-wise benchmark.
\end{itemize}

The broader message is that speculative decoding exposes a distinct
prompt-controlled systems vulnerability. Lossless output semantics do not imply
cost robustness. A user-controlled prefix can leave the final answer intact
while forcing the verifier to do substantially more work.

\section{Related Work}

\paragraph{Speculative decoding and acceptance bottlenecks.}
Speculative decoding was introduced as a lossless way to accelerate
autoregressive generation by verifying draft tokens in parallel with the target
model \citep{leviathan2023speculative}. Subsequent work has improved the draft,
the verifier, or the scheduling policy. Online Speculative Decoding adapts the
draft to the observed query distribution to increase token acceptance
\citep{liu2023online}. Direct alignment methods train draft models to better
match chat-capable targets on plausible data distributions
\citep{goel2024direct}. AdaEDL predicts when low acceptance is likely and stops
drafting early \citep{agrawal2024adaedl}. Judge Decoding argues that model
alignment alone is insufficient and redesigns verification to improve acceptance
\citep{bachmann2025judge}. Alignment-augmented decoding uses training-free
alignment sampling and flexible verification \citep{wang2025alignment}, while
HSD pushes verification beyond token-wise acceptance by performing lossless
hierarchical verification \citep{zhou2026hsd}.

\paragraph{Why our setting is different.}
These papers all treat low acceptance as a performance obstacle to be fixed.
Our viewpoint is the inverse: low acceptance is the attack objective. Rather
than training a better draft or verifier, we ask how a user-controlled prompt
can maximize draft-target disagreement under the deployed system. To our
knowledge, prior speculative-decoding work does not explicitly optimize prompt
prefixes to minimize early acceptance probabilities.

\paragraph{Prompt-based inference-cost attacks.}
The closest attack-style prior art is prompt-based latency manipulation in
ordinary autoregressive decoding. Engorgio constructs prompts that suppress
\texttt{<EOS>} and induce abnormally long outputs, thereby increasing inference
cost and threatening service availability \citep{dong2024engorgio}. This is
related in spirit but differs in mechanism. Engorgio attacks output length in
standard autoregressive generation. Our setting attacks the verification rule in
speculative decoding by making draft proposals difficult to accept.

\paragraph{Adversarial prompt optimization and universal triggers.}
Outside speculative decoding, a large literature studies adversarial prompts,
universal triggers, and jailbreak suffixes. Universal adversarial triggers show
that short input-agnostic token sequences can induce systematic failures across
many examples \citep{wallace2019uat}. In aligned-chat settings, discrete and
hybrid search methods such as the universal suffix attack of
\citet{zou2023universal} and the hierarchical genetic search of
\citet{liu2023autodan} demonstrate that prompt optimization can reliably steer
model behavior toward harmful semantic outputs. Our problem is different in
both objective and mechanism. We do not aim to bypass safety alignment or
change the semantic content of the answer. Instead, we optimize prompts to
alter the \emph{systems behavior} of speculative verification itself: the
target output can remain correct while the verifier is forced to do
substantially more work.

\paragraph{Positioning.}
The attack studied here lies at the intersection of systems efficiency and
adversarial prompting. The speculative-decoding literature establishes that
throughput depends strongly on draft-target agreement, while the adversarial
prompt literature shows that short optimized triggers can reliably manipulate
model behavior. ADSD combines these perspectives: it turns verifier agreement
into the attack surface and develops a continuous-to-discrete optimization
framework for prompt-controlled inference-cost degradation.

\section{Theoretical Framework and Implemented Objective}

To systematically find adversarial prefixes, we rely on gradient-based
optimization over continuous token embeddings. However, optimizing a continuous
surrogate does not guarantee that its nearest discrete projection remains
adversarial. In this section, we formalize the optimization objective,
acknowledge the theoretical limits of unconditional projection, and establish
the precise margin conditions under which our continuous-to-discrete pipeline is
guaranteed to succeed.

\subsection{Setup and the Continuous Attack Space}

Let the attacker optimize a prefix of length $m$. A discrete prefix is a
sequence $z = (z_1, \dots, z_m) \in \mathcal{V}^m$, where $\mathcal{V}$ is the
vocabulary. Let $E \in \mathbb{R}^{|\mathcal{V}| \times d}$ denote the
embedding matrix. We define a continuous relaxation
$U = (u_1, \dots, u_m)$, where each $u_t \in \Delta^{|\mathcal{V}|}$ is a
probability simplex vector. The relaxed embedding at position $t$ is
\[
e_t(U) = E^\top u_t.
\]

Let $p_i^U$ and $q_i^U$ denote the target and draft next-token distributions at
speculative step $i$ under the relaxed prefix $U$. Drawing from the one-step
expected acceptance identity $\mathbb{E}[\alpha_i] = 1 - \tv(p_i, q_i)$, a
natural theoretical objective is to maximize early draft-target disagreement:
\[
C(U) = \sum_{i=1}^{\gamma} w_i \tv(p_i^U, q_i^U),
\]
where $w_i = \gamma - i + 1$ linearly weights earlier tokens to maximize
rejection pressure. The corresponding discrete objective is defined by
restriction:
\[
C_d(z) := C(E(z)),
\]
where $E(z)$ is the concatenated embedding sequence of the discrete tokens.

To recover a deployable attack, we define the token-wise nearest-neighbor
projection $\Pi(U) := (\Pi_1(U), \dots, \Pi_m(U))$, where
\[
\Pi_t(U) = \arg\min_{v \in \mathcal{V}} \|e_t(U) - E_v\|_2.
\]
We quantify the continuous-to-discrete shift via the projection distortion
$\rho(U)$:
\[
\rho(U) := \sum_{t=1}^{m} \|e_t(U) - E_{\Pi_t(U)}\|_2.
\]

\subsection{The Obstruction of Prompt Waywardness}

A critical vulnerability in gradient-based text attacks is the assumption that
continuous optimization smoothly maps back to discrete space. We first state a
negative result demonstrating that without structural assumptions, projection is
fundamentally unsafe.

\begin{proposition}[Prompt Waywardness]
\label{prop:waywardness}
For any finite token embedding space with at least two tokens and any
$\delta > 0$, there exists a smooth objective $C$ over the convex hull of token
embeddings and a relaxed optimizer $U^*$ such that $U^*$ is within distance
$\delta$ of a discrete token cell, yet the nearest-neighbor projection
$\Pi(U^*)$ is arbitrarily suboptimal among all discrete prefixes.
\end{proposition}

This proposition formalizes the empirical warning signs from prior literature:
an unconditional projection guarantee is impossible. To bind the discrete gap,
we must constrain the behavior of the specific model distributions.

\subsection{Bounding the Projection Gap}

To guarantee attack transferability from continuous to discrete space, we assume
the generative distributions exhibit Lipschitz continuity with respect to the
input embeddings.

\begin{assumption}[TV-Lipschitzness]
\label{assump:tv-lipschitz}
For each speculative step $i$, there exist constants $L_i^p, L_i^q$ such that
for any two relaxed prefixes $U, V$,
\[
\tv(p_i^U, p_i^V) \le L_i^p \|U - V\|_{2,1},
\]
\[
\tv(q_i^U, q_i^V) \le L_i^q \|U - V\|_{2,1},
\]
where
\[
\|U - V\|_{2,1} := \sum_{t=1}^{m} \|e_t(U) - e_t(V)\|_2.
\]
Let
\[
L_C := \sum_{i=1}^{\gamma} w_i (L_i^p + L_i^q)
\]
denote the aggregate Lipschitz constant of the weighted objective.
\end{assumption}

Under this assumption, we can strictly bound the quality degradation incurred
during projection.

\begin{theorem}[Relaxation Gap for Early Acceptance Collapse]
\label{thm:relax-gap}
Let $U^*$ be an $\varepsilon$-optimal solution of the relaxed problem, such
that
\[
C(U^*) \ge \sup_U C(U) - \varepsilon.
\]
Let
\[
z^* \in \arg\max_{z \in \mathcal{V}^m} C_d(z)
\]
be the optimal discrete prefix. Then the projected discrete attack satisfies
\[
C_d(\Pi(U^*)) \ge C_d(z^*) - \varepsilon - L_C \rho(U^*).
\]
\end{theorem}

\begin{proof}
By Assumption~\ref{assump:tv-lipschitz} and the triangle inequality, $C$ is
$L_C$-Lipschitz. Let
$\hat{z} = \Pi(U^*)$. Because the continuous embeddings of $\hat{z}$ are within
$\rho(U^*)$ of $U^*$, it follows that
\[
C_d(\hat{z}) = C(E(\hat{z})) \ge C(U^*) - L_C \rho(U^*).
\]
Since the optimal discrete prompt is a valid feasible solution in the relaxed
space,
\[
C(U^*) \ge C_d(z^*) - \varepsilon.
\]
Combining these inequalities yields the bound.
\end{proof}

Theorem~\ref{thm:relax-gap} establishes that the discrete attack gap is
controlled by only two
actionable quantities: the optimization error of the relaxed solver
($\varepsilon$) and the distance from the relaxed solution to the token
manifold ($\rho(U^*)$).

We can further characterize when the continuous relaxation exactly recovers the
globally optimal discrete attack. Let the discrete attack margin be defined as
\[
\Gamma := C_d(z^*) - \max_{z \neq z^*} C_d(z).
\]

\begin{corollary}[Exact Recovery under Discrete Margin]
\label{cor:exact-recovery}
Assume the optimal discrete adversarial prefix $z^*$ is unique
($\Gamma > 0$). Under Theorem~\ref{thm:relax-gap}, if the combined optimization
error and
projection distortion is strictly less than the discrete margin,
\[
\varepsilon + L_C \rho(U^*) < \Gamma,
\]
then the nearest-neighbor projection perfectly identifies the optimal discrete
attack:
\[
\Pi(U^*) = z^*.
\]
\end{corollary}

\subsection{Algorithm Design Justification}

The theoretical bounds above directly dictate our optimization strategy. To
maximize the discrete adversarial effect $C_d(\Pi(U))$, we cannot simply
maximize $C(U)$. We must explicitly minimize the projection gap
$L_C \rho(U)$. Consequently, any practical attack should optimize a penalized
continuous objective of the form
\[
\max_U \; C(U) - \lambda \rho(U) - \tau \sum_{t=1}^{m} H(u_t),
\]
where $\lambda \rho(U)$ acts as a manifold-attraction penalty to explicitly
bound the gap in Theorem~\ref{thm:relax-gap}, and the entropy regularization
sharpens the simplex vectors so the optimizer explores regions near discrete
token boundaries.

\subsection{Implemented Soft-Prefix Attack}

Our implementation uses the same relaxed search space, but instantiates the
continuous objective for the \emph{first verifier step}, because early rejection
creates the largest downstream slowdown in token-wise speculative decoding. Let
$p^U$ and $q^U$ denote the target and draft next-token distributions after the
relaxed prefix $U$. We define the implemented objective
\begin{align*}
J(U) ={}&
\underbrace{\sum_{y} q^U(y)\,\mathrm{softplus}\!\left(\log q^U(y)-\log p^U(y)\right)}_{\text{soft collapse}} \\
&+ \beta \,\tv(p^U, q^U)
+ \eta \,\mathrm{KL}(q^U \Vert p^U) \\
&- \tau \,\frac{1}{m}\sum_{t=1}^{m} H(u_t)
- \lambda \,\rho(U).
\end{align*}
The first term directly rewards the failure mode we seek: tokens on which the
draft is confident while the verifier assigns much lower probability. The TV
and reverse-KL terms enlarge the broader disagreement region around that token,
and the entropy term sharpens the relaxed simplex vectors so that projection is
stable. The projection penalty $\rho(U)$ is the theory-guided term from
Theorem~\ref{thm:relax-gap}.

We optimize $J(U)$ with Adam over per-position vocabulary logits. Each step
uses a straight-through argmax estimator: the forward pass uses the projected
hard token at each position, while the backward pass differentiates through the
softmax relaxation. This is important operationally. The optimizer therefore
sees the same discrete token manifold that deployment will use, instead of
optimizing a purely continuous embedding string that may not survive
projection.

While Theorem~\ref{thm:relax-gap} bounds the discrete gap using explicit
$L_2$ nearest-neighbor projection in the embedding space, evaluating pairwise
Euclidean distances across the full tokenizer vocabulary at every optimization
step is computationally prohibitive. In the implementation, we therefore
approximate the operational projection with a straight-through argmax estimator
over the relaxed probability simplex. The theoretical connection is preserved
by retaining the explicit $L_2$ projection-distortion penalty $\rho(U)$ in the
continuous objective, which penalizes the distance between the relaxed
soft-embeddings and the hard embeddings induced by the argmax-selected tokens.

Model selection is also verifier-aligned. Although optimization uses the smooth
objective $J(U)$, we keep the best candidate by its exact projected
first-token acceptance
\[
\alpha_1 = \min\!\left(1, \frac{p(y_1)}{q(y_1)}\right),
\]
computed after hard projection. In other words, the smooth objective drives the
search, but the retained attack is always evaluated in the exact discrete
speculative-decoding pipeline.

In the current experiments, the key hyperparameters are a five-token prefix,
Adam learning rate $0.3$, temperature $1.0$, entropy weight $0.01$, TV weight
$1.0$, reverse-KL weight $0.1$, and gradient clipping at $1.0$. The
projection-distortion weight $\lambda$ is calibrated automatically from the
initial ratio between the unpenalized objective and the initial projection
distortion, then swept on Palmetto through a multiplicative factor. This
calibration gives a stable way to compare attacks across different prefix
lengths without hand-retuning the objective scale each time. Because the token
embedding space is highly non-convex, initializing the relaxed prefix from a
uniform distribution is unstable. We therefore warm-start the optimizer's
initial logits from a manual adversarial seed and then search the full
continuous vocabulary space around that initialization. In practice, the search
often discovers subword fragments or visually unusual tokens that strictly
outperform the original seed.

\section{Evaluation}
\label{sec:eval}

\subsection{Experimental Setup}

We evaluate ADSD in the reproduced HSD workspace of
\citet{zhou2026hsd} using \texttt{Qwen2.5-0.5B-Instruct-GPTQ-Int8} as the
draft model and \texttt{Qwen2.5-14B-Instruct-GPTQ-Int8} as the target model
under token-wise speculative verification. Unless otherwise noted, the
benchmark is GSM8K with the reproduced Qwen chat template and a matched
five-token benign prefix budget. All end-to-end numbers are reported against
the same benign token-wise baseline.

Each run uses batch size one, stochastic decoding
(\texttt{do\_sample=True}), and a per-sample cap of
\texttt{max\_new\_tokens=512}. We measure systems impact with wall-clock
latency, aggregate acceptance, block efficiency, and decoding speed. Latency
is recorded directly around the speculative \texttt{generate(...)} call.
Aggregate acceptance is the fraction of drafted tokens that survive
verification, and block efficiency is the number of accepted drafted tokens per
target verifier step. These counter-based metrics are computed from the saved
speculative traces in the reproduced HSD evaluator and therefore expose the
behavior of the verification mechanism independently of incidental timing
noise. Evaluating at batch size one follows standard speculative-decoding
practice \citep{leviathan2023speculative,zhou2026hsd}: the object of study is
the per-sequence verifier mechanism itself rather than downstream serving-stack
scheduling.

Before the end-to-end benchmarks, the Palmetto search consistently finds
projected prefixes that collapse early verifier acceptance. In the single-prompt
setting, the optimized five-token prefix reduces exact projected
$\alpha_1$ from $0.1597$ to $0.0288$. In the calibration-batch universal search
with $K=64$, every calibration size $N\in\{1,2,4,8\}$ yields a projected
prefix with calibration-batch $\alpha_1=0$, and the $N=4$ prefix delivers the
strongest accuracy-preserving end-to-end attack.

\subsection{End-to-End Latency Degradation}

ADSD substantially degrades verifier efficiency on unseen prompts while
preserving task accuracy. Table~\ref{tab:tokenwise-benchmark} reports the
matched 263-sample GSM8K benchmark comparing the benign boundary baseline with
the strongest accuracy-preserving universal prefix ($N=4$) and the strongest
cost-maximizing universal prefix ($N=2$).

\begin{table}[t]
\centering
\caption{Token-wise 263-sample GSM8K benchmark. Lower aggregate acceptance,
block efficiency, and decoding speed are worse for the system; higher mean
sample time is worse for the system.}
\label{tab:tokenwise-benchmark}
\footnotesize
\begin{tabular}{lccccc}
\toprule
Setting & Time (s) & Agg. acc. & Block eff. & Speed & GSM8K acc. \\
\midrule
Benign baseline & $23.10$ & $0.597$ & $5.832$ & $79.55$ & $222/263$ \\
Universal prefix ($N=4$) & $38.61$ & $0.390$ & $3.731$ & $51.68$ & $222/263$ \\
Universal prefix ($N=2$) & $51.60$ & $0.358$ & $3.506$ & $46.09$ & $202/263$ \\
\bottomrule
\end{tabular}
\end{table}

The systems signal is unambiguous. The $N=4$ universal prefix increases mean
sample time from $23.10$s to $38.61$s, a $67.1\%$ slowdown, while lowering
aggregate acceptance from $0.597$ to $0.390$ and block efficiency from $5.832$
to $3.731$. The same condition reduces HSD-style decoding speed from $79.55$
to $51.68$ while preserving GSM8K accuracy exactly at $222/263$. This
alignment between counter-based verifier statistics and wall-clock latency is
critical: it shows that the slowdown is a consequence of algorithmic verifier
degradation rather than an artifact of hardware noise. The $N=2$ universal
prefix pushes the cost attack further still, but at the price of task
degradation. For the core stealth claim of the paper, the $N=4$ prefix is the
most important result: a short prompt-side perturbation can deny the intended
speculative speedup while leaving user-visible task quality unchanged.

\subsection{Comparison with Discrete Search Baselines}

Bounded continuous optimization is strictly more effective than heuristic
discrete search for verifier manipulation. Table~\ref{tab:baseline-comparison}
compares ADSD against three discrete baselines: a pure random five-token
prefix, a GCG-adapted search baseline, and an AutoDAN-adapted search baseline.

\begin{table}[t]
\centering
\caption{Comparison with discrete search baselines on the token-wise GSM8K
benchmark. Search cost reports the dominant prompt-search workload required to
produce the final five-token prefix.}
\label{tab:baseline-comparison}
\footnotesize
\begin{tabular}{lccccc}
\toprule
Method & Search cost & Time (s) & Agg. acc. & Block eff. & GSM8K acc. \\
\midrule
Benign baseline & --- & $23.10$ & $0.597$ & $5.832$ & $222/263$ \\
Random prefix & $[XX]$ & $[XX]$ & $[XX]$ & $[XX]$ & $[XX]$ \\
GCG-adapted & $[XX]$ & $[XX]$ & $[XX]$ & $[XX]$ & $[XX]$ \\
AutoDAN-adapted & $[XX]$ & $[XX]$ & $[XX]$ & $[XX]$ & $[XX]$ \\
ADSD (UAT, $N=4$) & $[XX]$ & $38.61$ & $0.390$ & $3.731$ & $222/263$ \\
\bottomrule
\end{tabular}
\end{table}

The comparison isolates the key advantage of ADSD. Heuristic discrete methods
must search directly in token space, where verifier collapse is sparse and
projection failures are invisible until evaluation time. ADSD instead uses a
continuous verifier-aligned objective and a projection-aware regularizer to
reach deeper acceptance collapse with a short prompt budget. The resulting
prefixes therefore achieve stronger end-to-end degradation than discrete search
alone, while requiring materially less trial-and-error over the combinatorial
token space.

\subsection{Cross-Architecture Vulnerability}

The vulnerability is a property of speculative verification itself, not a quirk
of a single Qwen pairing. Table~\ref{tab:scaling-placeholder} evaluates ADSD on
larger Qwen and LLaMA draft-target pairs to test whether the same prompt-side
mechanism persists across architectures and scales.

\begin{table}[t]
\centering
\caption{Cross-architecture scaling of ADSD. Larger target models amplify the
absolute systems cost of verifier collapse because each rejected speculative
window forces more expensive target-side work.}
\label{tab:scaling-placeholder}
\footnotesize
\begin{tabular}{lcccc}
\toprule
Target $\rightarrow$ Draft & Condition & Time (s) & Block eff. & Accuracy \\
\midrule
Qwen2.5 14B $\rightarrow$ 0.5B & Benign & $23.10$ & $5.832$ & $222/263$ \\
Qwen2.5 14B $\rightarrow$ 0.5B & ADSD & $38.61$ & $3.731$ & $222/263$ \\
Qwen2.5 72B $\rightarrow$ 3B & Benign & $[XX]$ & $[XX]$ & $[XX]$ \\
Qwen2.5 72B $\rightarrow$ 3B & ADSD & $[XX]$ & $[XX]$ & $[XX]$ \\
LLaMA-3 8B $\rightarrow$ 1B & Benign & $[XX]$ & $[XX]$ & $[XX]$ \\
LLaMA-3 8B $\rightarrow$ 1B & ADSD & $[XX]$ & $[XX]$ & $[XX]$ \\
\bottomrule
\end{tabular}
\end{table}

The expected pattern is straightforward: if verifier collapse reflects a
fundamental weakness of draft-target acceleration rather than an idiosyncrasy of
one model family, then the same prompt-side mechanism should transfer across
architectures and become even more costly as target-model capacity grows. In
that regime, the absolute latency penalty becomes larger even when the relative
drop in acceptance is comparable.

\subsection{Ablation of Optimization Components}

The projection-aware design of ADSD is empirically necessary. Table
\ref{tab:ablation-placeholder} ablates the main terms of the optimization
objective, with special emphasis on the projection penalty $\rho(U)$$,$ which
is the direct algorithmic instantiation of the theoretical bound in
Section~\ref{sec:method}.

\begin{table}[t]
\centering
\caption{Ablation of ADSD optimization components. The full objective combines
early collapse pressure, trajectory-level distributional mismatch, and
projection-aware regularization.}
\label{tab:ablation-placeholder}
\footnotesize
\begin{tabular}{lcccc}
\toprule
Objective variant & Projected first token & Time (s) & Block eff. & Accuracy \\
\midrule
Full ADSD objective & $[XX]$ & $38.61$ & $3.731$ & $222/263$ \\
w/o projection penalty $\rho$ & $[XX]$ & $[XX]$ & $[XX]$ & $[XX]$ \\
w/o TV term & $[XX]$ & $[XX]$ & $[XX]$ & $[XX]$ \\
w/o reverse KL & $[XX]$ & $[XX]$ & $[XX]$ & $[XX]$ \\
w/o entropy term & $[XX]$ & $[XX]$ & $[XX]$ & $[XX]$ \\
\bottomrule
\end{tabular}
\end{table}

This ablation directly tests the central methodological claim of the paper. If
the relaxation-to-text gap were merely a nuisance detail, then removing the
projection-aware term would not materially affect the final discrete attack. Our
theory predicts the opposite: without a mechanism that keeps the relaxed prefix
close to the discrete manifold, optimization can find a strong continuous
adversarial state whose projected text no longer preserves verifier collapse.
The ablation therefore functions as an empirical validation of the
prompt-waywardness argument in Section~\ref{sec:method}.

\subsection{Discussion: Robustness Against Adaptive Defenses}

One natural defense is to detect persistent acceptance collapse and halt
drafting adaptively, as in designs such as AdaEDL. Such a fallback is useful,
but it does not eliminate the vulnerability studied here. The attacker does
not need to force the system into a permanently degraded regime to succeed; it
is already sufficient to deny the request the intended speculative speedup.
Moreover, adaptive fallbacks require an observation window before the system
can decide that drafting is no longer beneficial. The attack therefore still
extracts a penalty during the detection phase, after which the request is
reduced to ordinary autoregressive decoding rather than accelerated
speculative decoding. From the adversary's perspective, that outcome remains a
successful worst-case degradation of the algorithmic mechanism.


\section{Conclusion}

We introduced ADSD, a prompt-side attack view of speculative decoding. The
mathematical core is simple but consequential: expected accepted prefix length
is controlled by prefix products of acceptance probabilities, and any fixed
amount of rejection pressure is more damaging when moved earlier. This yields a
clear optimization target for adversarial prompt search and explains why
first-token collapse is the right objective when the attacker wants to maximize
fallback to the target model. Our preliminary experiments already show strong
first-step degradation on a real draft-target pair. The next step is to run the
full gradient-based search, extend evaluation to block-wise and hierarchical
verifiers, and measure end-to-end latency degradation directly.

\appendix

\section{Implementation Notes}

Our local attack evaluation uses the same draft-target model pair as the HSD
reproduction workspace and the same Qwen chat template as the GSM8K evaluator.
The preliminary adversarial seeds were ranked by exact first-token acceptance
using white-box access to both models. A separate search utility in the
companion codebase evaluates prompt candidates and records the draft token,
target probability, and divergence statistics for each seed.


\bibliographystyle{plainnat}
\bibliography{references}

\end{document}
